% =====================================================================================
%  RESERVE — Advanced-strategy material (Strategies 5, 6, 7)
%
%  NOT PART OF THE THESIS BUILD. This file is never \input by thesis_draft.tex and does
%  not compile on its own. It is a verbatim archive.
%
%  WHY THIS FILE EXISTS
%  The thesis compares the four Minimum-tier strategies (0-4). Strategies 5-7 -- residual
%  reinforcement learning on a cascaded baseline, learning-based model predictive control,
%  and residual reinforcement learning on the geometric baseline -- were drafted in full
%  before the scope was fixed at four. That prose is kept here, unedited, so the work is
%  not lost if the scope ever widens again.
%
%  Preserved 2026-09-14, from the last build that contained it (44 pages).
%
%  HOW TO PUT IT BACK
%  Each block below records the file and the section it came from. The prose refers to
%  labels that still exist in the live chapters (eq:ctrl-ades, sec:ctrl-s2, sec:rw-eval,
%  ...), so re-inserting a block at its original location should resolve without edits.
%  Three things WOULD need adjusting on re-insertion:
%    1. The strategy count in ch1 (Outline), ch4 (opening + summary) and ch5 (opening),
%       which currently reads "Strategies 0-4" / "Four Minimum-tier strategies".
%    2. Table 4.1 in ch4 -- add back the three rows recorded under BLOCK D.
%    3. The short scoping passages now standing in ch2 sec:rw-inject and sec:rw-rl, which
%       summarise these families in a few lines and say they are out of scope; those would
%       be replaced by the full text in BLOCKS A and B, not kept alongside it.
% =====================================================================================


% -------------------------------------------------------------------------------------
% BLOCK A --- ch2_related_work.tex, \section{Injecting the prediction: the controller
%             matters} (sec:rw-inject), the \paragraph{Inside a model predictive control
%             horizon.} paragraph, in full. Currently replaced by a ~10-line summary.
% -------------------------------------------------------------------------------------

\paragraph{Inside a model predictive control horizon.}
The residual can instead form part of the prediction model that an optimiser reasons over, so that
the controller plans against a disturbance it expects to meet rather than cancelling the one
present now. Chee et al.~\cite{chee2024knodedwmpc} embed a mixed-expert KNODE-DW model --- first
principles plus a deliberately lightweight neural ODE component --- inside the prediction
constraint of a nonlinear MPC, rather than adding it as an instantaneous feedforward term, and
evaluate it in both simulation and hardware. The hardware experiments use two Crazyflie~2.1
vehicles in stacked formation, the lower one carrying a thrust-upgrade bundle, flown with average
separations below \num{1.5} body lengths --- one case averaging \(\Delta z = \SI{0.12}{\metre}\),
about \num{1.2} body lengths --- with commanded separations from \SIrange{0.2}{0.4}{\metre} down to
a very tight \SI{0.08}{\metre}. Against a nominal MPC baseline they report average RMSE better by
\SI{40.1}{\percent} and maximum \(z\)-error smaller by \SI{57.5}{\percent}, from \SI{46}{\second}
of training data in total (\SI{18}{\second} simulated and \SI{28}{\second} on hardware). As with
the flatness-based work above, the controller runs \emph{off-board} --- an Intel~i5 laptop
commanding over Crazyradio at about \SI{400}{\hertz} with Vicon at about \SI{120}{\hertz}, the
onboard firmware retaining PID control of thrust and body rates.
Li et al.~\cite{li2023ndpnmpc} take a related route, training a multilayer perceptron --- spectrally
normalised on its weight matrices each epoch --- to predict the downwash force from the neighbour's
relative position and velocity, then injecting the predicted force sequence into the nonlinear MPC's
discrete prediction model as optimiser parameters, so that the lower vehicle can prepare in advance
from the neighbour's planned trajectory rather than from its current state alone. Their evaluation
is on hardware with two similar custom quadrotors, the optimisation running onboard a Jetson~TX2~NX
with pose from off-board motion capture. Hsieh et
al.~\cite{hsieh2025onlineadaptation} extend the approach by inserting an $\mathcal{L}_1$
residual-acceleration estimate into the KNODE-DW prediction model, and evaluate it on \emph{three}
Crazyflie~2.1 vehicles in V-stack and I-stack formations --- including a vertically aligned stack
with pairwise separations of \SI{0.2}{\metre}, about two body-lengths. Only the $\mathcal{L}_1$
estimate adapts online; the neural component reuses the training data of the work above and is not
retrained in flight. The evaluation is hardware-only, again off-board, with the vehicles driven
from separate machines and commanded in thrust and body rates. They report that \emph{``the
$\mathcal{L}_1$ adaptive module compensates for unmodeled disturbances most effectively when paired
with an accurate dynamics model''} --- in our reading, that online adaptation complements a good
residual model rather than substituting for one.


% -------------------------------------------------------------------------------------
% BLOCK B --- ch2_related_work.tex, \section{Learning the compensation directly}
%             (sec:rw-rl), the body of the section in full. The section itself still
%             exists in the live chapter, carrying a ~12-line summary instead.
% -------------------------------------------------------------------------------------

A third stance declines to model the force at all and learns the corrective action instead. Zhang
et al.~\cite{zhang2024proxfly} train a residual reinforcement-learning module on top of a cascaded
controller, producing high-level commands that compensate the disturbance and thrust loss caused by
downwash, and use domain randomisation to relax the requirement for accurate system identification.
They train in simulation and evaluate in both simulation and hardware, on a heterogeneous pair of
regular-sized quadcopters rather than the Crazyflie class. The residual and the high-level
controller run \emph{off-board} at \SI{50}{\hertz}, with the low-level loop onboard at
\SI{500}{\hertz} and motion capture at \SI{200}{\hertz}, so this is not an end-to-end onboard
residual controller.

One property distinguishes this approach sharply from everything in
Section~\ref{sec:rw-predictive}: the policy's observation comprises the tracking error, the
baseline cascaded controller's command and its own previous residual action, no neighbour state,
and no communication between vehicles~\cite{zhang2024proxfly}. Every relative-state method depends on knowing where the
neighbours are, whether by communication or by onboard sensing. Removing that dependency removes a
bandwidth requirement and a failure mode, at the cost of the ability to anticipate a neighbour the
vehicle cannot perceive.

The approach also gives up the explicit force estimate that the other families produce. That
estimate is what makes a predicted residual comparable against a measured one, and its absence
makes failures correspondingly harder to attribute, a consideration in experimental design as
much as in deployment.


% -------------------------------------------------------------------------------------
% BLOCK C --- ch4_control_architectures.tex, the three strategy sections in full,
%             between \section{Strategy 4...} and \section{What makes the comparison
%             fair...}. Nothing stands in their place in the live chapter; the chapter
%             runs Strategy 0-4 then straight to the fairness section.
% -------------------------------------------------------------------------------------

% -------------------------------------------------------------------------------------
\section{Strategy 5: residual reinforcement learning on a cascaded baseline}
\label{sec:ctrl-s5}

\paragraph{Implemented.} \textbf{Not implemented, and deliberately scheduled last.}

\paragraph{Reference.} Zhang et al.~\cite{zhang2024proxfly} train a residual reinforcement-learning
module on top of a \emph{cascaded} controller, not a geometric one. It produces corrections that
compensate the disturbance and thrust loss caused by downwash, and uses domain randomisation to
reduce the dependence on accurate system identification. Their residual and high-level controller
run off-board at \SI{50}{\hertz} above an onboard low-level loop at \SI{500}{\hertz}, on
regular-sized quadcopters. Strategy~5 follows that arrangement.

% -------------------------------------------------------------------------------------
\section{Strategy 6: learning-based model predictive control}
\label{sec:ctrl-s6}

\paragraph{Implemented.} \textbf{Not implemented.} Whether it can run within the control period on
this platform is an open engineering question, and one that is itself relevant. This is the only
strategy in the comparison whose computational cost may exclude it from the target hardware. Nothing
in the literature settles it either way for our vehicle: Chee et al.\ and Hsieh et al.\ both run
their optimisation off-board, and Li et al.\ run theirs on a Jetson companion computer, not on an
STM32.

\paragraph{Reference.} Chee et al.~\cite{chee2024knodedwmpc} embed a mixed-expert KNODE-DW model
--- first principles plus a lightweight neural ODE component --- inside the prediction constraint
of a nonlinear MPC, and evaluate it in simulation and on hardware with two Crazyflie~2.1 vehicles
in stacked formation at average separations below \num{1.5} body lengths, with the controller
running off-board. Hsieh et
al.~\cite{hsieh2025onlineadaptation} extend this by inserting an $\mathcal{L}_1$
residual-acceleration estimate into the same prediction model, evaluated on hardware with
\emph{three} Crazyflie~2.1 vehicles in V-stack and I-stack formations at pairwise separations down
to \SI{0.2}{\metre}, and report that the adaptive module is most effective when paired with an
accurate dynamics model. Only the $\mathcal{L}_1$ estimate adapts online.

\paragraph{Where the prediction enters.} Not at an instant, as in Strategies~2 and~4, but inside
the prediction model over which the optimiser reasons. The controller therefore plans against a
disturbance it expects to meet rather than cancelling the one present now, the most distinct
injection point in the comparison, and the reason the strategy is retained in the set even though
it may prove infeasible onboard.

% -------------------------------------------------------------------------------------
\section{Strategy 7: residual reinforcement learning on the geometric baseline}
\label{sec:ctrl-s7}

\paragraph{Implemented.} \textbf{Not implemented.} This pairing is ours, not any paper's.

\paragraph{What it shares with Strategy 5.} The stance is identical. A policy outputs a corrective
action, and no explicit force estimate is formed at any point. Both are therefore evaluated on
tracking performance alone, and neither can be checked by comparing a predicted residual against a
measured one.

\paragraph{What differs.} The inner loop. Strategy~5 sits on the cascaded controller of
\cite{zhang2024proxfly}; Strategy~7 sits on the geometric position and attitude loop used by
Strategies~0, 1, 2 and~4 in this thesis. Separating them matters because a residual policy is
trained against whatever the baseline leaves uncorrected, so the same learning stance on a
different inner loop is a different experiment, not a re-run.

\paragraph{Why they are scheduled last, stated plainly.} They answer a different question.
Strategies~1, 2, 3, 4 and~6 all reason about an explicit force: each produces, or consumes, an
estimate of $f_{\mathrm{res}}$ that can be compared against the measurement of
\eqref{eq:model-ares}. A residual policy produces a corrective action and never forms such an
estimate, so the predicted-versus-measured comparison that underpins the evaluation of every other
learned strategy is unavailable, and failures are correspondingly harder to attribute.

One property nonetheless makes them worth retaining. The policy of Zhang et al.\ observes only the
tracking error, the baseline controller's command and its own previous residual action, no
neighbour state, and no communication between
vehicles~\cite{zhang2024proxfly}. Every other predictive strategy here depends on knowing where the
neighbours are. Removing that dependency removes a bandwidth requirement and a failure mode, at the
cost of being unable to anticipate a neighbour the vehicle cannot perceive, a trade-off that
belongs in the cost accounting even if the strategies are not flown.


% -------------------------------------------------------------------------------------
% BLOCK D --- ch4_control_architectures.tex, Table 4.1 (tab:ctrl-summary): the three
%             rows removed from the strategy table. Re-insert after the Strategy 4 row.
% -------------------------------------------------------------------------------------

    5 & Residual RL, cascaded        & neither    & actuation                  & no  & deferred \\
    6 & Learning-based MPC           & predicted  & MPC prediction model       & no  & not implemented \\
    7 & Residual RL, geometric       & neither    & actuation                  & no  & deferred \\


% -------------------------------------------------------------------------------------
% BLOCK E --- ch4_control_architectures.tex, \section{Summary} (sec:ctrl-summary), the
%             seven-strategy wording. The live chapter carries the four-strategy version.
% -------------------------------------------------------------------------------------

Seven strategies, each its own controller, distinguished by what they know about the interaction
force, when they know it, and where that knowledge enters. Two of them, pure INDI and the hybrid,
happen to share an attitude law and differ only in which residual term is active --- a useful point
of comparison, not a privileged evidentiary category. All seven are compared under the same
experimental conditions: identical trajectories, formations, hardware and frozen gains, not
identical control structure.

At the time of writing, two strategies fly, one is implemented and unflown, one exists as separable
components, and three are not implemented.


% -------------------------------------------------------------------------------------
% BLOCK F --- one-line references removed from ch1, ch3 and ch5. Each is given with the
%             sentence it sat in, so it can be matched back to its location.
% -------------------------------------------------------------------------------------

% ch1_introduction.tex, \section{Outline}, the Chapter 3 entry:
%   "...consume a force residual use that number; Strategies~5 and~7 do not form it."
%   (live version stops after "use that number")

% ch1_introduction.tex, \section{Outline}, the Chapter 4 entry:
%   "describes Strategies~0--7 (eight controllers: an uncompensated baseline plus seven
%    families, with residual reinforcement learning split into two variants by inner loop)"
%   (live version: "describes Strategies~0--4 (the uncompensated baseline, pure INDI,
%    geometric+NN, FBL+NN and the hybrid)")

% ch3_system_modelling.tex, chapter opening:
%   "Residual reinforcement learning (Strategies~5 and~7) never forms that number. That is a
%    property of the method, not a reason to exclude it from a comparison of performance."

% ch3_system_modelling.tex, \section{Summary}:
%   "Strategies~5 and~7 never form it, which is a property of residual reinforcement learning
%    rather than a reason to exclude it from a comparison of performance."

% ch5_experimental_setup.tex, chapter opening:
%   "described Strategies~0--7 (eight controllers: an uncompensated baseline plus seven
%    families, with residual reinforcement learning split into two variants by inner loop)"
%   (live version: "described Strategies~0--4 (the uncompensated baseline, pure INDI,
%    geometric+NN, FBL+NN and the hybrid)")
